% =============================================================================
%  CHAPTER 19 -- PHASE 6: STATISTICAL DIAGNOSTICS ON THE TRADED RESIDUALS
% =============================================================================
\section{Phase 6 --- Statistical diagnostics on the traded residuals}
\label{ch:phase6}

\textbf{Reproduce it:} \code{python3 scripts/analysis/run\_diagnostics.py \&\& python3
scripts/plotting/render\_diagnostics.py}.\\
\textbf{Math used:} all of Chapters~\ref{ch:timeseries} and \ref{ch:ou}.\\
\textbf{Artifacts:} \code{results/tables/diagnostics\_\{summary,coint,halflife,multtest\}.csv},
\code{results/tables/spread\_residuals.csv},
\code{results/figures/diagnostics\_\{residual,halflife\}.png},
\code{report/tables/diagnostics.tex}.

\subsection{Goal}

Run the entire diagnostic battery on the residuals the strategy actually trades, and
report what it says even where it is uncomfortable. Four questions, in order of
importance:
\begin{enumerate}
  \item Is the spread stationary, or is the whole exercise a spurious regression
    (Section~\ref{sec:prob:spurious})?
  \item Are the target and basket genuinely cointegrated, and are the hedge weights
    significant under correct inference?
  \item How fast does the spread revert, and does that match the strategy's holding
    period?
  \item How much of the apparent significance disappears once autocorrelation and
    heteroskedasticity are accounted for?
\end{enumerate}

\subsection{Design: two residuals, two periods}

Every test is run on \emph{two} different objects, and the distinction is the most
important methodological point of the phase:

\begin{description}[leftmargin=0pt,style=nextline]
  \item[Static hedge] one OLS (and one Ridge) fit of the log-price relation over the whole
    period --- the object that ``is ADBE cointegrated with its peers?'' refers to, and the
    residual that Engle--Granger and Johansen are designed for.
  \item[Rolling causal spread] the 120-bar walk-forward hedge residual
    $S_t$ of \eqref{eq:saa:spread} --- the object the strategy \emph{actually trades},
    written to \code{results/tables/spread\_residuals.csv} bar by bar so the diagnostics
    and the backtest are provably about the same series.
\end{description}

Both are tested over the full sample (2010--2026, $n=4{,}070$ usable bars after warm-up)
and over the out-of-sample test period (2022+, $n=1{,}169$). Reporting only one of the
two would have produced a completely different --- and misleading --- story.

\subsection{Results: the stationarity battery}

\input{tables/diagnostics.tex}

\begin{figure}[H]\centering
\includegraphics[width=0.95\linewidth]{figures/diagnostics_residual.png}
\caption{The causal rolling spread actually traded (top) and its autocorrelation function
(bottom). The ACF decays toward the $\pm1.96/\sqrt T$ band by roughly lag 10--20,
consistent with the estimated OU half-life of $\approx9.4$ bars
(equation~\eqref{eq:ar1}: an AR(1)-type decay with $\phi=e^{-\theta}=0.928$ has
$\rho(10)=0.47$, $\rho(20)=0.22$).}
\label{fig:diag_residual}
\end{figure}

The measured verdicts (Table~\ref{tab:diag_stationarity}):

\begin{enumerate}[label=\textbf{D\arabic*.}]
  \item \textbf{The static long-run hedge is \emph{not} a clean cointegrating residual.}
    Over 2010--2026 the fixed OLS hedge residual gives ADF $-2.651$ with $p=0.083$ (fails
    to reject the unit root), KPSS $1.357$ with $p=0.01$ (rejects level-stationarity) ---
    both tests pointing the same way --- and an implied OU half-life of $170.8$ bars. On
    2022+ it is worse: ADF $p=0.845$, half-life $475.5$ bars. Ridge is nearly identical
    ($p=0.093$, half-life $178.1$ bars full sample). A single hedge ratio estimated once
    over sixteen years cannot pin Adobe to its peers; the relationship drifts.
  \item \textbf{The causal rolling spread \emph{is} mean-reverting.} ADF $-9.948$,
    $p<10^{-3}$ (full sample) and $-5.971$, $p<10^{-3}$ (2022+) --- decisive rejections of
    the unit root, with an OU half-life of $9.36$ bars full sample and $9.44$ bars on
    2022+, i.e.\ stable across regimes. KPSS, however, rejects level-stationarity on the
    full sample ($p=0.01$) and does not reject on 2022+ ($p=0.10$), so the joint verdict
    is \emph{ambiguous} full sample and \emph{stationary} on the test period. The
    ambiguity is informative rather than inconvenient: it is the signature of a process
    that reverts quickly toward a local mean that itself moves --- exactly what a rolling
    hedge produces, and exactly why the $z$-score of \eqref{eq:saa:z} uses a trailing
    $\hat\mu_{S,t}$ rather than a full-sample one.
  \item \textbf{Extreme serial correlation and strong heteroskedasticity, everywhere.}
    Durbin--Watson $=0.009$ (static) and $0.147$ (rolling), i.e.\ implied
    $\hat\rho_1\approx0.995$ and $0.93$ by \eqref{eq:dw:rho}; Ljung--Box $Q=38{,}379$ and
    $20{,}725$ with $p\approx0$; Breusch--Pagan and White both $p<10^{-3}$. For the
    static residual, DW$\approx0.01$ is the textbook signature of a \emph{unit-root}
    residual (successive values nearly identical). For the rolling residual it is the
    signature of the persistent mispricing we intend to trade. Both are the same
    statistic; the interpretation depends entirely on which object you are looking at ---
    which is why the two-residual design matters.
  \item \textbf{The consequence for inference is a factor of 3--4 in $t$-statistics.}
    Table~\ref{tab:diag_coint}: the classic (i.i.d.) $t$ for the \code{CRM} hedge weight
    is $29.52$ and the Newey--West HAC $t$ is $7.29$; for \code{ADSK}, $29.69\to9.42$; for
    the intercept, $-44.95\to-12.27$. Naive inference overstates significance by
    $\approx4\times$ in $t$ and $\approx16\times$ in variance terms --- the
    $1+2\sum_\ell\rho(\ell)$ inflation factor of \eqref{eq:longrunvar} with
    $\rho(\ell)\approx0.99^\ell$. Every significance claim in this report is made on HAC
    standard errors, never on the classic ones.
\end{enumerate}

\subsection{Results: cointegration and hedge significance}

\begin{itemize}
  \item \textbf{Engle--Granger, name by name} (target \code{ADBE} against each peer
    individually): \code{CRM} $p=0.0032$ and \code{ADSK} $p=0.0325$ are cointegrated at
    5\%; \code{INTU} $p=0.1165$ is not.
  \item \textbf{Johansen on the four-name log-price vector}
    $(\text{ADBE},\text{CRM},\text{ADSK},\text{INTU})$: cointegration rank
    $\mathbf{r=2}$ at the 95\% level (trace statistic, \eqref{eq:johansen:trace}). So the
    basket does share two independent stationary combinations --- it is genuinely
    cointegrated at the system level, not merely highly correlated.
  \item \textbf{Hedge weights under HAC inference} (Table~\ref{tab:diag_coint}):
    $\hat\beta_{\text{CRM}}=0.649$ ($t=7.29$), $\hat\beta_{\text{ADSK}}=0.600$
    ($t=9.42$), $\hat\beta_{\text{INTU}}=0.027$ ($t=0.45$), intercept $-0.898$
    ($t=-12.27$). The hedge is economically real but carried by two of the three names.
\end{itemize}

\begin{remark}[Honesty: rank 2 does not mean our spread is one of them]
The Johansen test says two stationary directions exist. It does not say that the
particular direction our rolling-OLS hedge constructs is one of them, and the
$\alpha$-loading half of the VECM (which would tell us how fast each name corrects) is not
estimated here. The honest summary is: the basket is cointegrated, the two dominant hedge
weights are individually significant under HAC inference, \code{INTU}'s weight is not,
and the traded spread is stationary by direct test --- which is the claim the strategy
needs. That is stronger than ``the names are correlated'' and weaker than ``we have
identified the cointegrating vector''.
\end{remark}

\begin{remark}[Honesty: \code{INTU} is insignificant but still useful]
$\hat\beta_{\text{INTU}}=0.027$ with $t=0.45$ looks like a name that should be dropped ---
and Phase~\ref{ch:phase5} found that Lasso/ElasticNet shrinkage in this direction
\emph{degraded} the book. The resolution is that \code{INTU} is individually
collinear with \code{CRM}/\code{ADSK} (VIF of Section~\ref{sec:ols:multicol} in the
tens), so its marginal contribution is small while its variance-reducing contribution to
the rolling hedge is not. This is a concrete case where a $t$-statistic is the wrong tool
for the decision --- a point in favour of ridge's ``keep everything, shrink smoothly''
prior over lasso's ``select'' prior for hedging.
\end{remark}

\subsection{Results: half-life versus holding period}

\begin{figure}[H]\centering
\includegraphics[width=0.85\linewidth]{figures/diagnostics_halflife.png}
\caption{OU half-life (log scale) versus realized average holding period, for the static
hedge and the traded rolling spread, full sample and 2022+. The two static-hedge bars
($171$--$476$ bars) are not tradeable quantities; the rolling-spread half-life
($\approx9.4$ bars) is, and the realized hold of $16$--$18$ bars is $1.7$--$1.9$ of it.}
\label{fig:diag_halflife}
\end{figure}

Measured: traded (rolling) spread $\hat\tau_{1/2}=9.36$ bars (full) / $9.44$ bars (2022+);
realized average holding period $16.2$ bars (full) / $17.3$ bars (2022+). By
\eqref{eq:ou:gain} the strategy therefore captures $1-e^{-\theta h}=1-2^{-1.73}\approx70\%$
(full sample) to $1-2^{-1.85}\approx72\%$ (test) of the expected reversion per trade. That
is a sensible horizon by the theory of Chapter~\ref{ch:ou} --- not so short that it trades
noise, not so long that it pays carry for a reversion that has already happened --- but it
is on the turnover-heavy side given the cost drag measured in Phase~\ref{ch:phase4}, and
the sensitivity analysis of Phase~\ref{ch:phase7} confirms that moving to fewer, longer
trades is the biggest available net-Sharpe lever.

The static hedge's nominal half-life of $\sim$171--476 bars is reported only to make the
contrast vivid; it is not a strategy parameter, because a half-life estimated on a
non-stationary residual measures the drift of the hedge, not the speed of a reversion.

\subsection{A first multiple-testing note}

This phase also recorded a preliminary hypothesis count
(\code{results/tables/diagnostics\_multtest.csv}): 13 candidates scanned in the
Phase~\ref{ch:phase5} walk-forward grid on one basket at fixed cost, with a Bonferroni
threshold of $0.05/13=0.00385$. Phase~\ref{ch:phase10} supersedes it with the full audit
over $M=44$ strategies on a common calendar, including White's Reality Check and Hansen's
SPA. The interim number is kept in the results tree because it documents that the count
was being tracked from the first phase in which a choice was made, not retrofitted at the
end.

\subsection{Exit gate --- status}

\begin{itemize}
  \item[\checkmark] all specified diagnostics run and tabulated for the final models
    (stationarity, cointegration, autocorrelation, heteroskedasticity, HAC, OU);
  \item[\checkmark] HAC before/after $t$-statistics reported, not just the corrected ones;
  \item[\checkmark] half-life versus holding-period mismatch explicitly discussed;
  \item[\checkmark] a numeric multiple-testing threshold recorded (superseded by
    Phase~\ref{ch:phase10});
  \item[\checkmark] the diagnostics operate on the \emph{same} series the backtest trades
    (via \code{spread\_residuals.csv}).
\end{itemize}
